\section*{Заключение}
\addcontentsline{toc}{section}{Заключение}

В ходе производственной практики разработан полный ML-пайплайн для задачи ранжирования коротких видео: от индексации исходных данных до обучения ранжирующей модели с трекингом экспериментов.

\textbf{В части инженерии данных} выполнено следующее:
\begin{itemize}
    \item реализован этап индексации данных с сортировкой недельных шардов по глобальному временнóму порядку;
    \item настроена потоковая очистка 48~млн событий на Polars (аномалии времени просмотра, фильтры минимальной активности) без загрузки данных в память;
    \item построена градуированная целевая переменная из неявного фидбека с настраиваемыми весами сигналов;
    \item рассчитаны три группы признаков (пользовательские, айтемные, парные) строго по train-окну, без утечки будущего; признак близости контентного эмбеддинга видео к эмбеддинг-профилю пользователя вычисляется на PyTorch двумя потоковыми проходами;
    \item выполнено темпоральное разбиение 80/10/10 и собраны датасеты для CatBoost (38.4~млн / 4.8~млн / 4.8~млн событий, 36 числовых и 5 категориальных признаков).
\end{itemize}

\textbf{В части машинного обучения}:
\begin{itemize}
    \item обучена модель \texttt{CatBoostRanker} с listwise-лоссом YetiRank (группы --- пользователи), поддерживающая обучение на GPU с автоматическим переходом на CPU;
    \item реализован собственный векторизованный модуль оффлайн-метрик ранжирования (NDCG@k, MAP@k, MRR@k, HitRate@k, Recall@k, Precision@k, GAUC), сверенный с эталонной реализацией;
    \item метрики рассчитываются как в процессе обучения (по итерациям, на валидации), так и после него на валидации и тесте; кривые и финальные значения логируются в ClearML;
    \item итоговая модель, обученная на GPU-сервере, достигла на тестовом сплите NDCG@10~=~0.408, MRR@10~=~0.804, GAUC~=~0.807 при согласованности метрик валидации и теста.
\end{itemize}

\textbf{В части инфраструктуры и качества}:
\begin{itemize}
    \item локально развернут сервер ClearML в Docker Compose (семь сервисов, включая эмуляцию amd64 на arm64-машине); каждый этап пайплайна фиксируется как отдельный эксперимент;
    \item пайплайн и модель покрыты 40 автоматизированными тестами на синтетических данных с подсаженным сигналом, включая проверки отсутствия утечек таргета и будущего, поэлементную сверку формулы релевантности и smoke-обучение с проверкой превосходства модели над случайным ранжированием;
    \item вся конфигурация вынесена в YAML-файлы; проверочные прогоны управляются параметрами \texttt{max\_shards}, \texttt{user\_fraction}, \texttt{max\_rows} без изменения кода.
\end{itemize}

Таким образом, все планируемые результаты практики достигнуты: настроен воспроизводимый пайплайн обработки больших данных, подготовлены очищенные датасеты для ML-задач, обучена и оценена ранжирующая модель, решение задокументировано.

Направлениями дальнейшего развития могут быть:
\begin{enumerate}
    \item обучение на полном датасете с распределенной обработкой признаков;
    \item расширение признакового описания: временные признаки (динамика интереса, свежесть видео), последовательные модели истории просмотров;
    \item подбор гиперпараметров через ClearML HyperParameter Optimization;
    \item двухстадийная архитектура (отбор кандидатов + ранжирование) и оценка на A/B-симуляции;
    \item автоматизация запуска этапов через ClearML Pipelines с удаленными агентами.
\end{enumerate}
